\subsection{商作用}

定义4. 设 \(\alpha \mapsto f_{\alpha}\) 是集合 \(\Omega\) 在集合 \(\mathbf{E}\) 上的一个作用。一个等价关系 \(\mathbf{R}\) （在 \(\mathbf{E}\) 上）称为与给定作用相容，如果对所有元素 \(x\) 和 \(y\) （属于 \(\mathbf{E}\) ，使得 \(x \equiv y \pmod{\mathbf{R}}\) ）以及所有 \(\alpha \in \Omega\) , \(f_{\alpha}(x) \equiv f_{\alpha}(y) \pmod{\mathbf{R}}\) 。将元素 \(\alpha \in \Omega\) 与 \(\mathbf{E} / \mathbf{R}\) 到自身的、由 \(f_{\alpha}\) 通过取商导出的映射相关联的映射是 \(\Omega\) 在 \(\mathbf{E} / \mathbf{R}\) 上的一个作用，称为 \(\Omega\) 在 \(\mathbf{E}\) 上的作用的商.

设 E 是一个广群，R 是 E 上的一个等价关系。若 R 与由 E 上的律导出的 E 在自身上的左（相应地，

（右）作用相容，则称 R 与 E 上的律左（相应地，右）相容。R 与 E 上的律相容当且仅当它与 E 上的律既左相容又右相容。

我们将 §1 第 6 号中命题 6、7 和 8 的类似陈述与证明留给读者。
% semantic-fragments: legacy-input-seam
\relax{}
